\chapter{Заключение}
\label{ch:final}

За учебный год в ходе выполнения лабораторных работ я освоил следующие навыки и технологии.

\section*{Android-разработка}
\begin{itemize}
    \item Работа с геолокацией: получение координат (широта, долгота, высота), точности, текущего времени с помощью FusedLocationProviderClient.
    \item Использование разрешений ACCESS\_FINE\_LOCATION, ACCESS\_COARSE\_LOCATION, READ\_PHONE\_STATE.
    \item Сбор информации о сотовых сетях через TelephonyManager: данные LTE, GSM, NR.
    \item Создание фонового сервиса для сбора данных в фоне, с запуском и остановкой из Activity.
    \item Передача данных между сервисом и Activity через LocalBroadcastManager.
    \item Сохранение данных в JSON-файлы во внутреннее хранилище.
    \item Подключение к ZMQ-серверу и отправка JSON-строк.
    \item Разработка MediaPlayer: воспроизведение mp3 из внешней памяти, управление громкостью, SeekBar, список треков.
    \item Создание нескольких Activities и переходы между ними, обработка нажатий кнопок.
\end{itemize}

\section*{Backend-разработка}
\begin{itemize}
    \item Написание сервера на C++ с использованием библиотеки ZeroMQ.
    \item Парсинг входящих JSON-сообщений с помощью библиотеки nlohmann/json.
    \item Многопоточность: отдельный поток для ZMQ-сервера и отдельный поток для графического интерфейса, синхронизация через мьютексы и общую структуру LocationData.
    \item Сохранение полученных данных в JSON-файл с накоплением истории.
    \item Работа с базой данных PostgreSQL: подключение через libpq, создание таблиц based и cells, вставка измерений и данных о сотах.
    \item Визуализация графиков сигналов для разных PCI с помощью библиотеки ImPlot, цветовая дифференциация.
    \item Отображение карт OpenStreetMap: асинхронная загрузка тайлов через libcurl, декодирование PNG в RGBA, создание текстур OpenGL.
    \item Реализация динамической подгрузки тайлов при изменении масштаба и панорамировании, кэширование на диске.
    \item Генерация тепловых карт методом обратных взвешенных расстояний по критериям RSRP, RSRQ, RSSI, высоте. Наложение полупрозрачной тепловой карты поверх OSM.
    \item Настройка параметров интерполяции.
\end{itemize}

\section*{Инструменты}
\begin{itemize}
    \item Работа с системой контроля версий Git: создание веток, коммиты, слияния, push в удалённый репозиторий.
    \item Ведение README.md с описанием проекта, инструкциями по сборке и запуску.
    \item Использование библиотек: ZeroMQ, nlohmann/json, libpq, libcurl, libpng, SDL2, ImGui, ImPlot.
    \item Взаимодействие Android-приложения с сервером через TCP/IP в локальной сети.
\end{itemize}

Таким образом, я приобрёл практические навыки разработки полноценного программно-аппаратного комплекса для сбора, передачи, хранения и визуализации данных мобильного телефона, что позволяет в дальнейшем применять эти знания для реальных задач мониторинга качества сотовой связи.
\endinput